\section{Theory}\label{sec:theory}

This section answers the question of the paper in closed form for the predictor of Section~\ref{sec:method} with its
scalars given. Theorem~\ref{thm:risk} gives its expected squared error at any operating point, exactly, for any design and
any plant; Corollaries~\ref{cor:align} to~\ref{cor:coverage} read it in three limits; Theorem~\ref{thm:fixed} averages it over
the design and separates what the engine's curvature buys from what its plane buys; Proposition~\ref{prop:line} sets the
method against the setting-only line; Proposition~\ref{prop:priors} relates the three discrepancy priors; and
Proposition~\ref{prop:bound} states what a design can and cannot identify. Every statement holds conditionally on the
training design, and on the scalars except in Propositions~\ref{prop:priors} and~\ref{prop:bound}, which concern the
linearized model of Section~\ref{sec:method-linear} in which the scalars are learned; the full-sampler runs of
Section~\ref{sec:results-sampler} measure what the linearization leaves out.

\subsection{The risk of the predictor given its scalars}

Fix one output and drop its index. The plant has a mean function $\mu(\mathbf x)$ and returns $y=\mu(\mathbf x)+\varepsilon$
with $\mathbb E\varepsilon=0$ and $\mathrm{Var}\,\varepsilon=\sigma^2$, independently over surveys. The training design
$\mathbf x_1,\dots,\mathbf x_n$ gives the regressor matrix $\mathbf Z$ of~\eqref{eq:design}, with $\mathbf G=\mathbf
Z^{\top}\mathbf Z$ nonsingular, and the engine at the given scalars gives $f(\mathbf x)$. Project both functions on the
column space of $\mathbf Z$ over the training surveys:
\begin{equation}
\begin{gathered}
\boldsymbol\beta=\mathbf G^{-1}\mathbf Z^{\top}\boldsymbol\mu,\quad \mathbf r=\boldsymbol\mu-\mathbf Z\boldsymbol\beta,\\
\mathbf b=\mathbf G^{-1}\mathbf Z^{\top}\mathbf f,\quad \mathbf e=\mathbf f-\mathbf Z\mathbf b,
\end{gathered}
\label{eq:projections}
\end{equation}
so that $\mathbf Z^{\top}\mathbf r=\mathbf Z^{\top}\mathbf e=\mathbf 0$. At a held-out point $\mathbf x_h$ with design row
$\tilde{\mathbf z}_h$ write $r_h=\mu(\mathbf x_h)-\tilde{\mathbf z}_h^{\top}\boldsymbol\beta$ for the plant's departure from its
training plane, $e_h=f(\mathbf x_h)-\tilde{\mathbf z}_h^{\top}\mathbf b$ for the engine's, and
$\boldsymbol\delta=\mathbf b-\boldsymbol\beta$ for the difference of the two planes. The vector $\mathbf r$ is the part of the
plant that no regression on the setting and the feed size can represent; $\mathbf e$ is the part of the engine that the
discrepancy plane cannot represent, its curvature in the two regressors; $\boldsymbol\delta$ is how far the engine's best
plane lies from the plant's. With $\mathbf A_\lambda=(\mathbf G+\lambda\mathbf I)^{-1}$ the predictor~\eqref{eq:cbmean} is
$\hat y_\lambda(\mathbf x_h)=f(\mathbf x_h)+\tilde{\mathbf z}_h^{\top}\mathbf A_\lambda\mathbf Z^{\top}(\mathbf y-\mathbf f)$ and
least squares on $\mathbf Z$ is $\hat y_{\mathrm{LS}}(\mathbf x_h)=\tilde{\mathbf z}_h^{\top}\mathbf G^{-1}\mathbf Z^{\top}\mathbf y$.

\begin{theorem}[Exact risk given the scalars]\label{thm:risk}
For every $\lambda\ge0$, with $w_h=\tilde{\mathbf z}_h^{\top}\mathbf A_\lambda\boldsymbol\delta$ and the expectation over the
training and the held-out noise,
\begin{align}
R_\lambda(\mathbf x_h)&=\big(r_h-e_h-\lambda w_h\big)^2+\sigma^2\,\tilde{\mathbf z}_h^{\top}\mathbf A_\lambda\mathbf G\mathbf
A_\lambda\tilde{\mathbf z}_h+\sigma^2,\label{eq:riskcb}\\
R_{\mathrm{LS}}(\mathbf x_h)&=r_h^2+\sigma^2\,\tilde{\mathbf z}_h^{\top}\mathbf G^{-1}\tilde{\mathbf z}_h+\sigma^2.\label{eq:riskls}
\end{align}
If the noise is Gaussian, the error $y_h-\hat y_\lambda(\mathbf x_h)$ is Gaussian with mean $m_h=r_h-e_h-\lambda w_h$ and
variance $v_h=\sigma^2(1+\tilde{\mathbf z}_h^{\top}\mathbf A_\lambda\mathbf G\mathbf A_\lambda\tilde{\mathbf z}_h)$.
\end{theorem}
\begin{IEEEproof}
Write $\boldsymbol\mu=\mathbf Z\boldsymbol\beta+\mathbf r$ and $\mathbf f=\mathbf Z\mathbf b+\mathbf e$. Since $\mathbf
Z^{\top}\mathbf r=\mathbf Z^{\top}\mathbf e=\mathbf 0$, $\mathbf Z^{\top}(\mathbf y-\mathbf f)=-\mathbf G\boldsymbol\delta+\mathbf
Z^{\top}\boldsymbol\varepsilon$. With $\mathbf A_\lambda\mathbf G=\mathbf I-\lambda\mathbf A_\lambda$,
\begin{align*}
\hat y_\lambda(\mathbf x_h)&=\tilde{\mathbf z}_h^{\top}\mathbf b+e_h-\tilde{\mathbf z}_h^{\top}\boldsymbol\delta+\lambda w_h
+\tilde{\mathbf z}_h^{\top}\mathbf A_\lambda\mathbf Z^{\top}\boldsymbol\varepsilon\\
&=\tilde{\mathbf z}_h^{\top}\boldsymbol\beta+e_h+\lambda w_h+\tilde{\mathbf z}_h^{\top}\mathbf A_\lambda\mathbf Z^{\top}\boldsymbol\varepsilon .
\end{align*}
The held-out survey is $y_h=\tilde{\mathbf z}_h^{\top}\boldsymbol\beta+r_h+\varepsilon_h$, so the error is
$r_h-e_h-\lambda w_h-\tilde{\mathbf z}_h^{\top}\mathbf A_\lambda\mathbf Z^{\top}\boldsymbol\varepsilon+\varepsilon_h$: a constant
plus two independent zero-mean terms with variances $\sigma^2\tilde{\mathbf z}_h^{\top}\mathbf A_\lambda\mathbf Z^{\top}\mathbf
Z\mathbf A_\lambda\tilde{\mathbf z}_h$ and $\sigma^2$, which gives~\eqref{eq:riskcb} and, under Gaussian noise, the stated law.
Least squares is the same computation with the engine removed, $\hat y_{\mathrm{LS}}=\tilde{\mathbf z}_h^{\top}\boldsymbol\beta+\tilde{\mathbf
z}_h^{\top}\mathbf G^{-1}\mathbf Z^{\top}\boldsymbol\varepsilon$, which gives~\eqref{eq:riskls}.
\end{IEEEproof}

The theorem splits the risk into three terms with distinct causes. The first is a squared bias: the plant's curvature
$r_h$ that neither predictor can fit, corrected by the engine's own curvature $e_h$, and moved by the shrinkage toward the
engine's plane through $\lambda w_h$. The second is the variance of the fitted plane, which the shrinkage reduces. The third
is the noise of the held-out survey, which no predictor removes. Nothing in the theorem assumes that the plant is the engine,
that the discrepancy prior is right or that the noise scale of the model equals the plant's: $\lambda$ is any non-negative
number, and the plant is any function. The proof also shows that the predictor with $\lambda=0$ is least squares on
$\mathbf Z$ plus $e_h$ exactly, survey by survey and noise draw by noise draw.

\begin{corollary}[The engine helps only through aligned curvature]\label{cor:align}
At $\lambda=0$, $R_{\mathrm{LS}}(\mathbf x_h)-R_0(\mathbf x_h)=e_h(2r_h-e_h)$. For held-out points drawn from a distribution
$P$ with $\|r\|_P>0$ and $\|e\|_P>0$, with $\langle u,v\rangle_P=\mathbb E_P[uv]$, $\rho=\langle e,r\rangle_P/(\|e\|_P\|r\|_P)$ and
$\kappa=\|e\|_P/\|r\|_P$,
\begin{equation}
\mathbb E_P\big[R_{\mathrm{LS}}-R_0\big]=2\langle e,r\rangle_P-\|e\|_P^2=\|r\|_P^2\,\big(2\rho\kappa-\kappa^2\big),
\label{eq:align}
\end{equation}
which is positive if and only if $0<\kappa<2\rho$, and largest, $\rho^2\|r\|_P^2$, at $\kappa=\rho$. When the plant is a
plane on $P$ ($\|r\|_P=0$) the difference is $-\|e\|_P^2\le0$.
\end{corollary}
\begin{IEEEproof}
The variance terms of~\eqref{eq:riskcb} at $\lambda=0$ and~\eqref{eq:riskls} coincide, and $r_h^2-(r_h-e_h)^2=e_h(2r_h-e_h)$;
take the expectation and divide by $\|r\|_P^2$. The quadratic $2\rho\kappa-\kappa^2$ in $\kappa$ has roots $0$ and $2\rho$.
\end{IEEEproof}

The corollary is the answer to the question of when the physics pays for itself, with the shrinkage switched off. The
engine improves on an empirical regression in the same regressors only through the part of it that is not a plane, and only
when that part points the same way as the plant's own departure from a plane ($\rho>0$) and is not more than twice as large
as the aligned part of the plant's ($\kappa<2\rho$). An engine whose curvature has the right sign but twice the right size
gains nothing; one whose curvature has the wrong sign always loses. The gain is not a variance term: it depends on $n$ only
through the training planes $\boldsymbol\beta$ and $\mathbf b$, and settles to its population value as the design fills its
range, instead of decaying like $1/n$ as the variance of a fitted plane does. At a small $n$ it varies with the design and can
take either sign. That persistence is the property that makes a physics-informed predictor worth its cost.
The noise does not enter~\eqref{eq:align} at all. What enters is the plant's mean function, which no finite set of surveys
reveals; this is why the question cannot be settled on ten surveys, and why the study of Section~\ref{sec:design} uses
virtual plants whose mean function is known.

\begin{corollary}[An affine engine is only a shrinkage target]\label{cor:affine}
If the engine is affine in the regressors at every operating point, $e\equiv0$, then
\begin{equation}
\begin{split}
R_\lambda(\mathbf x_h)-R_{\mathrm{LS}}(\mathbf x_h)&=\lambda^2w_h^2-2\lambda r_hw_h\\
&\quad-\sigma^2\,\tilde{\mathbf z}_h^{\top}\big(\mathbf G^{-1}-\mathbf A_\lambda\mathbf G\mathbf A_\lambda\big)\tilde{\mathbf z}_h ,
\end{split}
\label{eq:affine}
\end{equation}
and $\mathbf G^{-1}-\mathbf A_\lambda\mathbf G\mathbf A_\lambda$ is positive semidefinite, with eigenvalues
$(2\lambda g_k+\lambda^2)/(g_k(g_k+\lambda)^2)$ for the eigenvalues $g_k$ of $\mathbf G$.
\end{corollary}
\begin{IEEEproof}
Set $e_h=0$ in~\eqref{eq:riskcb} and subtract~\eqref{eq:riskls}. $\mathbf G$, $\mathbf A_\lambda$ and their products share the
eigenvectors of $\mathbf G$, on which $1/g-g/(g+\lambda)^2=(2\lambda g+\lambda^2)/(g(g+\lambda)^2)$.
\end{IEEEproof}

On the HP500 lane the engine's throughput is $Q_{\mathrm{ref}}(c+17)/38$ at the fixed speed, affine in the setting at every
value of the eight scalars (Proposition~2 of~\cite{cb002}), so the corollary applies to throughput exactly. CrusherBayes then
predicts throughput as a ridge regression toward the engine's plane: it can save part of the variance of the fitted plane,
and the engine contributes no curvature of its own. The only change of the bias is the shrinkage term $\lambda w_h$, which
averages to zero over the training points (Theorem~\ref{thm:fixed}) and elsewhere, under extrapolation in particular, moves
the bias by the mismatch of the two planes, in a direction that no property of the engine controls. A capacity law that is curved in the
setting, as the power law of Rocha et al.~\cite{rocha}, is the condition under which Corollary~\ref{cor:align} can act on
throughput at all; Section~\ref{sec:results-capacity} measures how much it acts.

\begin{corollary}[The coverage of the band given the scalars]\label{cor:coverage}
Under Gaussian noise, the band $\hat y_\lambda(\mathbf x_h)\pm q_h$ covers the held-out survey with probability
\begin{equation}
\Pr\big(|y_h-\hat y_\lambda(\mathbf x_h)|\le q_h\big)=\Phi\!\left(\frac{q_h-m_h}{\sqrt{v_h}}\right)+\Phi\!\left(\frac{q_h+m_h}{\sqrt{v_h}}\right)-1,
\label{eq:coverage}
\end{equation}
with $m_h$ and $v_h$ of Theorem~\ref{thm:risk} and $\Phi$ the standard normal distribution function. With the predictive
variance of the model at the plant's noise scale, $s_h^2=\sigma^2(1+\tilde{\mathbf z}_h^{\top}\mathbf A_\lambda\tilde{\mathbf
z}_h)$, one has $s_h^2-v_h=\sigma^2\lambda\,\tilde{\mathbf z}_h^{\top}\mathbf A_\lambda^2\tilde{\mathbf z}_h\ge0$.
\end{corollary}
\begin{IEEEproof}
The first statement is the law of Theorem~\ref{thm:risk}. For the second, $\mathbf A_\lambda-\mathbf A_\lambda\mathbf G\mathbf
A_\lambda=\mathbf A_\lambda(\mathbf G+\lambda\mathbf I-\mathbf G)\mathbf A_\lambda=\lambda\mathbf A_\lambda^2$.
\end{IEEEproof}

The band of the model therefore errs in two directions at once: it is wider than the spread of the error by the prior's
share $\sigma^2\lambda\tilde{\mathbf z}_h^{\top}\mathbf A_\lambda^2\tilde{\mathbf z}_h$, and it is centred on a biased mean
whenever $m_h\ne0$. The first effect raises the coverage above the nominal level and the second lowers it, so a band that
covers 80\pct{} of a set of held-out surveys can do so for the wrong reason. When the plant is the engine plus a plane drawn
from the prior of the model ($\mathbf r=\mathbf e$) and the model's noise scale is the plant's, averaging over the drawn
planes makes the bias and the extra width cancel and the coverage nominal; at a fixed discrepancy, as on a real plant, or with
a plant whose curvature differs from the engine's, they do not, and~\eqref{eq:coverage} gives the coverage the band actually has. The study reports both the coverage and the
bias-eliminated coverage of Morris et al.~\cite{morris}, which separates the two effects.

\subsection{The risk averaged over the design}

Theorem~\ref{thm:risk} holds point by point. To see how the terms scale with the number of surveys, average it over the
training surveys themselves as held-out points with new noise, the fixed-design risk $\bar R=\frac1n\sum_iR(\mathbf x_i)$ of
regression theory~\cite{hastie}. Let $\mathbf G=\mathbf V\mathrm{diag}(g_1,g_2,g_3)\mathbf V^{\top}$ and
$\tilde{\boldsymbol\delta}=\mathbf V^{\top}\boldsymbol\delta$.

\begin{theorem}[Curvature gain and plane gain]\label{thm:fixed}
With $C=\frac1n\sum_{i=1}^n e_i(2r_i-e_i)$,
\begin{equation}
\bar R_\lambda-\bar R_{\mathrm{LS}}=-C+\frac1n\sum_{k=1}^{3}\frac{\lambda^2g_k\tilde\delta_k^2-\sigma^2(2\lambda g_k+\lambda^2)}{(g_k+\lambda)^2}.
\label{eq:fixed}
\end{equation}
For $\tilde\delta_k\ne0$ the $k$-th term of the sum is minimized at $\lambda_k^\star=\sigma^2/\tilde\delta_k^2$, where it equals
$-\sigma^4/(n(g_k\tilde\delta_k^2+\sigma^2))$; for $\tilde\delta_k=0$ it decreases toward $-\sigma^2/n$ as $\lambda$ grows. The
whole sum is never below $-3\sigma^2/n$.
\end{theorem}
\begin{IEEEproof}
Average~\eqref{eq:riskcb} over the training points. The cross term $-\frac{2\lambda}{n}(\mathbf r-\mathbf e)^{\top}\mathbf
Z\mathbf A_\lambda\boldsymbol\delta$ vanishes because $\mathbf Z^{\top}(\mathbf r-\mathbf e)=\mathbf 0$; the squared shrinkage
term is $\frac{\lambda^2}{n}\boldsymbol\delta^{\top}\mathbf A_\lambda\mathbf G\mathbf A_\lambda\boldsymbol\delta=\frac{\lambda^2}{n}\sum_kg_k\tilde\delta_k^2/(g_k+\lambda)^2$;
the variance term is $\frac{\sigma^2}{n}\mathrm{tr}(\mathbf A_\lambda\mathbf G\mathbf A_\lambda\mathbf G)=\frac{\sigma^2}{n}\sum_kg_k^2/(g_k+\lambda)^2$;
and $\frac1n\|\mathbf r-\mathbf e\|^2=\frac1n\|\mathbf r\|^2-C$. Least squares gives $\frac1n\|\mathbf r\|^2+3\sigma^2/n+\sigma^2$.
Subtracting and writing $\sigma^2g_k^2-\sigma^2(g_k+\lambda)^2=-\sigma^2(2\lambda g_k+\lambda^2)$ gives~\eqref{eq:fixed}. For one
term, $\varphi(\lambda)=(\lambda^2g\tilde\delta^2+\sigma^2g^2)/(g+\lambda)^2$ has $\varphi'(\lambda)\propto\lambda g\tilde\delta^2-\sigma^2g$,
so its minimum is at $\lambda^\star=\sigma^2/\tilde\delta^2$ with $\varphi(\lambda^\star)=\sigma^2g\tilde\delta^2/(g\tilde\delta^2+\sigma^2)$;
subtracting $\sigma^2$ gives the stated value, which lies between $-\sigma^2$ and $0$.
\end{IEEEproof}

Theorem~\ref{thm:fixed} separates the two ways the engine can help. The curvature gain $C$ is the empirical form of
Corollary~\ref{cor:align} and settles, as the design grows, to its population value; it does not decay with $n$. The plane
gain is a variance saving bounded by the whole variance $3\sigma^2/n$ of least squares on three regressors; with a design
whose second moments settle, $g_k$ grows like $n$, and when $\tilde\delta_k\ne0$ the gain in that direction decays like
$\sigma^4/(n^2\nu_k\tilde\delta_k^2)$, with $\nu_k=\lim g_k/n$. With few surveys the plane can matter; with many, only the
curvature can.

The theorem also says how the discrepancy prior should be scaled. The prior of Kennedy and O'Hagan, $\mathbf d\sim
N(\mathbf 0,\tau^2\mathbf I)$, with the noise scale of the model equal to the plant's, gives $\lambda=\sigma^2/\tau^2$, which is
the shrinkage that minimizes the fixed-design risk in the eigendirection $k$ of $\mathbf G$ exactly when
\begin{equation}
\tau^2=\tilde\delta_k^2 .
\label{eq:taustar}
\end{equation}
The risk-optimal discrepancy scale is the mismatch between the engine's plane and the plant's plane, measured in the
eigendirections of $\mathbf G$, for the fixed-design risk; at held-out points the cross term of Corollary~\ref{cor:affine}
enters and the optimum moves. A single $\tau$ cannot match three different mismatches, and the theorem gives the
loss of every compromise. This is a frequentist reading of a Bayesian prior: the prior is right, for prediction, when its
scale is the size of the error it is meant to absorb.

\subsection{Against the setting-only line}

The line in the setting alone is the comparator that the held-out surveys could not separate from CrusherBayes. Let
$\mathbf Z_1$ hold the first two columns of $\mathbf Z$, $\mathbf P$ and $\mathbf P_1$ the projections on the column spaces of
$\mathbf Z$ and $\mathbf Z_1$, and $\hat y_{\mathrm{line}}$ least squares on $\mathbf Z_1$.

\begin{proposition}[When the third regressor and the engine pay]\label{prop:line}
With $\phi^2=\frac1n\|(\mathbf P-\mathbf P_1)\boldsymbol\mu\|^2$, the mean square of the part of the plant that the feed size
explains beyond the setting,
\begin{equation}
\bar R_{\mathrm{line}}-\bar R_\lambda=\phi^2-\frac{\sigma^2}{n}+C-\frac1n\sum_{k=1}^{3}\frac{\lambda^2g_k\tilde\delta_k^2-\sigma^2(2\lambda g_k+\lambda^2)}{(g_k+\lambda)^2}.
\label{eq:line}
\end{equation}
At $\lambda=0$ and at a given design the method beats the line if and only if $n(\phi^2+C)>\sigma^2$. For a design whose
empirical distribution settles, with $\phi^2$ and $C$ at their limiting values, this reads $n>n^\star=\sigma^2/(\phi^2+C)$ when
$\phi^2+C>0$, and never when $\phi^2+C\le0$.
\end{proposition}
\begin{IEEEproof}
Least squares on $\mathbf Z_1$ has $\bar R_{\mathrm{line}}=\frac1n\|\boldsymbol\mu-\mathbf P_1\boldsymbol\mu\|^2+2\sigma^2/n+\sigma^2$, and
since $\mathbf P_1\mathbf P=\mathbf P_1$, $\|\boldsymbol\mu-\mathbf P_1\boldsymbol\mu\|^2=\|\mathbf r\|^2+\|(\mathbf P-\mathbf
P_1)\boldsymbol\mu\|^2$. Hence $\bar R_{\mathrm{line}}-\bar R_{\mathrm{LS}}=\phi^2-\sigma^2/n$, and~\eqref{eq:fixed} gives the rest.
\end{IEEEproof}

The proposition puts a number on the observation that ten surveys could not separate the method from the line. The third
regressor costs one parameter's variance, $\sigma^2/n$, and earns $\phi^2$; the engine earns its curvature gain $C$; the
shrinkage earns at most the variance of the plane. On power the noise of the surveys is large against both $\phi^2$ and
$C$ (Section~\ref{sec:results-theory} computes them on the virtual plant fitted to the ten surveys), so the crossing size
$n^\star$ lies beyond ten. The proposition is an equality at every design, so with the limiting values of $\phi^2$ and $C$ it
also says how many surveys would be needed.

\subsection{Three discrepancy priors}

The predictor~\eqref{eq:cbmean} takes the scalars as given. When they are learned, CrusherBayes in its linearized form
(Section~\ref{sec:method-linear}) is a Gaussian linear model in $\mathbf w=(\mathbf u,\mathbf d)$ with features $\mathbf
W=[\mathbf J,\mathbf Z]$ and prior precision $\mathbf K$, and its predictive mean is
\begin{equation}
m_{\mathbf K}(\mathbf x_h)=f_c(\mathbf x_h)+\mathbf w_h^{\top}\big(\mathbf W^{\top}\mathbf W+\sigma^2\mathbf K\big)^{-1}\mathbf W^{\top}(\mathbf y-\mathbf f_c),
\label{eq:genridge}
\end{equation}
a generalized ridge regression on the engine's sensitivities and the discrepancy regressors, toward the engine at the
linearization point $\boldsymbol\theta_c$. The three priors of Section~\ref{sec:method-priors} differ only in $\mathbf K$ and in
the features.

\begin{proposition}[What the priors change]\label{prop:priors}
(a) The orthogonal prior, with features $[\mathbf J,(\mathbf I-\mathbf J\mathbf J^{+})\mathbf Z]$ and held-out row $(\mathbf
j_h,\tilde{\mathbf z}_h-(\mathbf J^{+}\mathbf Z)^{\top}\mathbf j_h)$ under precision $\mathbf K_\perp$, gives the predictive mean
of the Kennedy and O'Hagan features $\mathbf W$ under the precision $\mathbf T^{-\top}\mathbf K_\perp\mathbf T^{-1}$, with
$\mathbf T=\begin{pmatrix}\mathbf I&-\mathbf J^{+}\mathbf Z\\ \mathbf 0&\mathbf I\end{pmatrix}$.
(b) If the null space $\mathcal N$ of $\mathbf W$ satisfies $\mathbf w(\mathbf x)\perp\mathcal N$ at every operating point
$\mathbf x$ (a structural confounding), then as $\mathbf K=\epsilon\mathbf K_0$ with $\epsilon\downarrow0$ the predictive
mean~\eqref{eq:genridge} tends to the same limit for every $\mathbf K_0\succ0$; every prior predicts alike in the diffuse limit,
however it splits the confounded parameters.
(c) With no engine parameter learned, the scaled prior $\mathbf K_d=\tau^{-2}(\mathbf I+c\,\mathbf G/n)$ is the predictor of
Theorem~\ref{thm:fixed} with the shrinkage $\lambda_k=\lambda(1+c\,g_k/n)$ in the $k$-th eigendirection of $\mathbf G$,
and~\eqref{eq:fixed} holds term by term with $\lambda_k$ in place of $\lambda$.
\end{proposition}
\begin{IEEEproof}
(a) $[\mathbf J,(\mathbf I-\mathbf J\mathbf J^{+})\mathbf Z]=\mathbf W\mathbf T$ and the held-out row is $\mathbf T^{\top}\mathbf w_h$;
inserting in~\eqref{eq:genridge}, $\mathbf w_h^{\top}\mathbf T(\mathbf T^{\top}\mathbf W^{\top}\mathbf W\mathbf T+\sigma^2\mathbf
K_\perp)^{-1}\mathbf T^{\top}\mathbf W^{\top}=\mathbf w_h^{\top}(\mathbf W^{\top}\mathbf W+\sigma^2\mathbf T^{-\top}\mathbf K_\perp\mathbf
T^{-1})^{-1}\mathbf W^{\top}$, and $\mathbf T$ is invertible. (b) As $\epsilon\downarrow0$, $(\mathbf W^{\top}\mathbf W+\epsilon\sigma^2\mathbf
K_0)^{-1}\mathbf W^{\top}\mathbf v$ tends to the solution of the normal equations $\mathbf W^{\top}\mathbf W\mathbf w=\mathbf W^{\top}\mathbf
v$ of least $\mathbf K_0$-norm; two such solutions differ by an element of $\mathcal N$, which $\mathbf w(\mathbf x)^{\top}$ maps to zero.
(c) With $\mathbf W=\mathbf Z$, $\mathbf W^{\top}\mathbf W+\sigma^2\mathbf K_d=\mathbf G+\lambda(\mathbf I+c\mathbf G/n)$, diagonal in the
eigenvectors of $\mathbf G$ with entries $g_k+\lambda_k$.
\end{IEEEproof}

Three confoundings of CrusherBayes are structural in the sense of (b) in the joint model of the two outputs. First, the
engine computes power as $\zeta W_i\,10(P_{80}^{-1/2}-F_{80}^{-1/2})\,Q_{\mathrm{ref}}(c+17)/38+P_n$, so the capacity base,
the work index and the power slope enter power only through their product $Q_{\mathrm{ref}}W_i\zeta$ (the product of
Proposition~4 of~\cite{cb002}, extended to the capacity base), while the sensitivity of throughput to $\ln Q_{\mathrm{ref}}$,
$Q_{\mathrm{ref}}\tilde{\mathbf z}^{\top}\mathbf a$, lies in the span of the throughput discrepancy at every operating point;
a move of $+1$ in $\ln Q_{\mathrm{ref}}$, $-1$ in $\ln W_i$ and $-Q_{\mathrm{ref}}\mathbf a$ in the throughput discrepancy is
therefore invisible to both outputs everywhere. Second, $+1$ in $\ln W_i$ and $-1$ in $\ln\zeta$ is invisible everywhere.
Third, the no-load power has the sensitivity $1$, equal to the constant regressor of the power discrepancy, at every
operating point. With at least as many surveys as the features have identified combinations (eight on power, which every
design of the study exceeds), the null space of $\mathbf W$ contains only these structural directions, and part (b) applies,
the capacity base included. The orthogonal prior of Plumlee~\cite{plumlee} therefore cannot change the predictions of the
method in the diffuse limit; its effect is on what the scalars mean, which part (a) makes explicit as a correlation between
the scalars and the discrepancy in the prior. With proper priors the three priors are three generalized ridges with different
penalties; by analogy with~\eqref{eq:taustar}, which is derived for the discrepancy alone, their predictive risks can be
expected to differ by how well each penalty matches the mismatch, a statement the study measures rather than proves. Part (c) states the scaled prior of
Gu and Wang~\cite{guwang} in this finite form: it shrinks most in the directions the design measures best, which is the
opposite of what~\eqref{eq:taustar} asks when the large-eigenvalue directions carry the large mismatches.

\subsection{What a design can identify}

Proposition~1 of~\cite{cb002} gives the Fisher information of a scalar under the marginal likelihood with the discrepancy
integrated out. Its consequence for design is a bound that does not depend on the number of surveys.

\begin{proposition}[An information ceiling]\label{prop:bound}
Let the sensitivity of output $t$ to a scalar $\vartheta$ lie in the span of its discrepancy regressors at every design point,
$\mathbf j_t=\mathbf Z_t\mathbf c_t$. Then its Fisher information about $\vartheta$ satisfies
\begin{equation}
\mathcal I_t(\vartheta)=\mathbf c_t^{\top}\mathbf G_t\big(\sigma_t^2\mathbf I+\tau_t^2\mathbf G_t\big)^{-1}\mathbf c_t\ \le\ \frac{\|\mathbf c_t\|^2}{\tau_t^2}
\label{eq:ceiling}
\end{equation}
for every design (with $\mathbf c_t$ that design's coefficients) and every number of surveys. In the linear Gaussian model the
posterior standard deviation of $\vartheta$ is then at least $(s_\vartheta^{-2}+\|\mathbf c_t\|^2/\tau_t^2+\mathcal
I_{\mathrm{other}})^{-1/2}$, with $s_\vartheta$ its prior standard deviation and $\mathcal I_{\mathrm{other}}$ the information
from the other outputs. A design factor that moves the sensitivity out of the span gives the undiscounted term
$\|(\mathbf I-\boldsymbol\Pi_t)\mathbf j_t\|^2/\sigma_t^2$, with $\boldsymbol\Pi_t$ the projection on the column space of
$\mathbf Z_t$, which grows in proportion to the number of surveys.
\end{proposition}
\begin{IEEEproof}
$\mathbf j_t^{\top}\boldsymbol\Sigma_t^{-1}\mathbf j_t$ with $\boldsymbol\Sigma_t=\sigma_t^2\mathbf I+\tau_t^2\mathbf Z_t\mathbf Z_t^{\top}$
and the push-through identity $\mathbf Z_t^{\top}\boldsymbol\Sigma_t^{-1}\mathbf Z_t=\mathbf G_t(\sigma_t^2\mathbf I+\tau_t^2\mathbf G_t)^{-1}$
give the equality; each eigenvalue $g/(\sigma_t^2+\tau_t^2g)$ is below $1/\tau_t^2$. In the linear Gaussian model the posterior
precision of $\vartheta$ is at most its prior precision plus its Fisher information, which gives the standard deviation. The
last statement is the first term of Proposition~1 of~\cite{cb002}.
\end{IEEEproof}

For the HP500 capacity base at the fixed speed, $\mathbf c=Q_{\mathrm{ref}}\mathbf a$ with $\mathbf a=((\bar c+17)/38,\varsigma_c/38,0)$
depends only on the mean $\bar c$ and the spread $\varsigma_c$ of the settings, so no number of surveys at the same speed and
in the same setting range can narrow $\ln Q_{\mathrm{ref}}$ from throughput beyond~\eqref{eq:ceiling}. The standard-deviation
bound of the proposition is valid but does not bind on its own, because $\mathcal I_{\mathrm{other}}$, the power's information,
grows with $n$. What binds is the joint structural confounding stated after Proposition~\ref{prop:priors}: power informs only
the sum $\ln Q_{\mathrm{ref}}+\ln W_i+\ln\zeta$. With that sum known exactly, the precision of $\ln Q_{\mathrm{ref}}$ is its
prior precision, the throughput ceiling, and the precision that the priors of $\ln W_i$ and $\ln\zeta$ lend to the
difference, so that for every number of surveys at the fixed speed, in the linear Gaussian model,
\begin{equation}
\frac{\mathrm{sd}(\ln Q_{\mathrm{ref}}\mid\mathcal D)}{s_Q}\ \ge\ \Big(1+s_Q^2\,\frac{\|\mathbf c\|^2}{\tau_{\mathrm{tph}}^2}+\frac{s_Q^2}{s_{W}^2+s_{\zeta}^2}\Big)^{-1/2},
\label{eq:floor}
\end{equation}
with $s_Q$, $s_W$ and $s_\zeta$ the prior standard deviations of the three logarithms.
Section~\ref{sec:results-design} evaluates the ceiling on the ten surveys and sets it beside the contraction the sampler
measured; it also gives the expected information gain of each candidate design. For a linear Gaussian model the expected
information gain of Rainforth et al.~\cite{rainforth}, the mutual information between the scalars and the data, has the
closed form
\begin{equation}
\mathrm{EIG}(\xi)=\tfrac12\ln\det\!\Big(\mathbf I+\mathbf L^{\top}\Big(\sum_t\mathbf J_t^{\top}\boldsymbol\Sigma_t^{-1}\mathbf J_t\Big)\mathbf L\Big),
\label{eq:eig}
\end{equation}
with $\mathbf L$ the diagonal of the prior standard deviations and $\mathbf J_t$, $\boldsymbol\Sigma_t$ evaluated on the design
$\xi$~\cite{chaloner}. It counts the information of every output, discounted by its discrepancy as
Proposition~1 of~\cite{cb002} prescribes, and it ranks designs that add an output (the product size), a factor (the speed)
or surveys.
